\section{Number Theory}
Used to speed up calculations on polynomials.

\subsection{Chinese Remainder Theorem}
Works the exact same way as for integers.

\subsection{CRT Polynomial}
Used to reduce the size of (the representation of) large polynomials, as $q$ is often much larger than a computer word (32- or more commonly 64-bit words)

\subsection{Double CRT Polynomial}
Used to perform multiplication on the smaller coefficient-first representation of large polynomials.

\subsection{Number Theoretic Transform}
The fast fourier transform (FTT) is an algorithm often used to simplify calculations in signals processing etc. by moving a continuous function from the time domain to the frequency domain. The \gls{ntt} works in a similar fashion; it's essentially a FFT mod q used to speed up multiplication of polynomials.

It is performed on DCRT polynomials.